Independence is what makes a differential oracle worth believing, and it is graded
along the wrong axis. Whether the reference was written independently bounds how
much logic it can inherit; what the reference reads bounds which defects it can
inherit, and that is the bound that decides what the oracle will never see. The two
are orthogonal, and only the first has a name.

We named the second, gave a procedure that turns it into a number before the oracle
is built, and applied it once: an oracle built with every care available was
measured, from source alone, to be structurally blind to 42\% of a documented defect
population and able to target 2\% of it---figures conditional on a frame we then
measured to be enriched for the oracle's own path, reproduced by a blind second rater,
and obtainable at the cost of opening fifty-seven commits. Reading gave the bound; explaining why one of its boundaries holds took a
build, and the paper does not claim the first act does the second's work.

The procedure's limit is also its finding. It cannot be validated beyond one system
today, because no defect population anywhere is organized by whether an oracle could
have observed each defect. Taxonomies record what a defect did, not what could have
caught it. The study that would close this is specific rather than gestural: take a
defect population drawn independently of any tool---one release window's
wrong-code defects in a compiler, say---and partition it by whether an oracle of a
given input provenance could have observed each. That measures the blind side, which
a tool's own findings can never do, and it is the first application of the same
procedure to a system its authors did not build.
